\section{Minimum Path Sum}
\label{sec:dp-minimum-path-sum}

\problemstatement{Given an $m \times n$ grid of non-negative numbers,
find a path from the top-left to the bottom-right corner, moving only
right or down, that \textbf{minimizes} the sum of the numbers along the
path.}

\begin{patternbox}
This is Unique Paths (Section~\ref{sec:dp-unique-paths}) with sum
replaced by minimum -- exactly the same transformation that turned
Climbing Stairs into Frog Jump back in the DP 1D chapter. The
recurrence's \emph{structure} (combine the cell above and the cell to
the left) is identical; only the combination operator and the addition
of the current cell's own value change.
\end{patternbox}

\subsection*{Deriving the recurrence}

Let $\textit{minSum}(i,j)$ be the minimum path sum to reach $(i,j)$. The
path's last step into $(i,j)$ came from $(i-1,j)$ or $(i,j-1)$, and we
must add the current cell's own value to whichever predecessor gives the
smaller total:
\[
  \textit{minSum}(i,j) = \textit{grid}[i][j] + \min\big(
  \textit{minSum}(i-1,j),\ \textit{minSum}(i,j-1)\big).
\]
Base case: $\textit{minSum}(0,0) = \textit{grid}[0][0]$. For cells along
the top row or left column, only one predecessor exists (the other is
out of bounds and treated as $+\infty$, never the chosen minimum).

\subsection*{Approach 1: Recursion}

\begin{lstlisting}
#include <bits/stdc++.h>
using namespace std;

int minPathSumRecursive(int i, int j, vector<vector<int>>& grid) {
    if (i == 0 && j == 0) return grid[0][0];
    if (i < 0 || j < 0) return INT_MAX;

    int fromAbove = minPathSumRecursive(i - 1, j, grid);
    int fromLeft  = minPathSumRecursive(i, j - 1, grid);
    return grid[i][j] + min(fromAbove, fromLeft);
}

int main() {
    vector<vector<int>> grid = {{1,3,1},{1,5,1},{4,2,1}};
    int m = grid.size(), n = grid[0].size();
    cout << minPathSumRecursive(m - 1, n - 1, grid) << endl; // 7
    return 0;
}
\end{lstlisting}

\begin{complexitybox}[Approach 1 Complexity]
\textbf{Time.} $O(2^{m+n})$.
\textbf{Space.} $O(m+n)$ recursion stack.
\end{complexitybox}

\subsection*{Approach 2: Memoization}

\begin{lstlisting}
#include <bits/stdc++.h>
using namespace std;

int minPathSumMemo(int i, int j, vector<vector<int>>& grid,
                    vector<vector<int>>& memo) {
    if (i == 0 && j == 0) return grid[0][0];
    if (i < 0 || j < 0) return INT_MAX;
    if (memo[i][j] != -1) return memo[i][j];

    int fromAbove = minPathSumMemo(i - 1, j, grid, memo);
    int fromLeft  = minPathSumMemo(i, j - 1, grid, memo);
    return memo[i][j] = grid[i][j] + min(fromAbove, fromLeft);
}

int main() {
    vector<vector<int>> grid = {{1,3,1},{1,5,1},{4,2,1}};
    int m = grid.size(), n = grid[0].size();
    vector<vector<int>> memo(m, vector<int>(n, -1));
    cout << minPathSumMemo(m - 1, n - 1, grid, memo) << endl; // 7
    return 0;
}
\end{lstlisting}

\begin{complexitybox}[Approach 2 Complexity]
\textbf{Time.} $O(m \cdot n)$.
\textbf{Space.} $O(m \cdot n)$ memo $+$ $O(m+n)$ recursion stack.
\end{complexitybox}

\subsection*{Approach 3: Tabulation}

\begin{lstlisting}
#include <bits/stdc++.h>
using namespace std;

int minPathSumTabulation(vector<vector<int>>& grid) {
    int m = grid.size(), n = grid[0].size();
    vector<vector<int>> dp(m, vector<int>(n, 0));

    for (int i = 0; i < m; i++) {
        for (int j = 0; j < n; j++) {
            if (i == 0 && j == 0) {
                dp[i][j] = grid[0][0];
                continue;
            }
            int fromAbove = (i > 0) ? dp[i-1][j] : INT_MAX;
            int fromLeft  = (j > 0) ? dp[i][j-1] : INT_MAX;
            dp[i][j] = grid[i][j] + min(fromAbove, fromLeft);
        }
    }
    return dp[m - 1][n - 1];
}

int main() {
    vector<vector<int>> grid = {{1,3,1},{1,5,1},{4,2,1}};
    cout << minPathSumTabulation(grid) << endl; // 7
    return 0;
}
\end{lstlisting}

\subsection*{Dry run of the tabulation table}

\dryrun{Take $\textit{grid} = \begin{bmatrix}1&3&1\\1&5&1\\4&2&1\end{bmatrix}$.}

\begin{center}
\begin{tabular}{c|c|c|c}
 & $j=0$ & $j=1$ & $j=2$ \\ \hline
$i=0$ & 1 & 4 & 5 \\ \hline
$i=1$ & 2 & 7 & 6 \\ \hline
$i=2$ & 6 & 8 & 7
\end{tabular}
\end{center}

E.g.\ $\textit{dp}[1][2] = \textit{grid}[1][2] + \min(\textit{dp}[0][2],
\textit{dp}[1][1]) = 1+\min(5,7) = 1+5=6$. The answer is
$\textit{dp}[2][2]=7$, achieved by the path $1\to3\to1\to1\to1$ (right,
right, down, down), summing to $1+3+1+1+1=7$.

\begin{complexitybox}[Approach 3 Complexity]
\textbf{Time.} $O(m \cdot n)$.
\textbf{Space.} $O(m \cdot n)$ table, no recursion stack.
\end{complexitybox}

\subsection*{Approach 4: Space Optimization}

\begin{lstlisting}
#include <bits/stdc++.h>
using namespace std;

int minPathSumSpaceOptimized(vector<vector<int>>& grid) {
    int m = grid.size(), n = grid[0].size();
    vector<int> prevRow(n, 0);

    for (int i = 0; i < m; i++) {
        vector<int> currentRow(n, 0);
        for (int j = 0; j < n; j++) {
            if (i == 0 && j == 0) {
                currentRow[j] = grid[0][0];
                continue;
            }
            int fromAbove = prevRow[j];
            int fromLeft  = (j > 0) ? currentRow[j-1] : INT_MAX;
            currentRow[j] = grid[i][j] + min(fromAbove, fromLeft);
        }
        prevRow = currentRow;
    }
    return prevRow[n - 1];
}

int main() {
    vector<vector<int>> grid = {{1,3,1},{1,5,1},{4,2,1}};
    cout << minPathSumSpaceOptimized(grid) << endl; // 7
    return 0;
}
\end{lstlisting}

\begin{complexitybox}[Approach 4 Complexity]
\textbf{Time.} $O(m \cdot n)$.
\textbf{Space.} $O(n)$ for the rolling row.
\end{complexitybox}

\begin{takeawaybox}
Once a grid-path-counting recurrence is in hand, converting it into a
grid-path-\emph{cost} recurrence is purely mechanical: replace the sum
of predecessors with a minimum (or maximum) of predecessors, plus the
current cell's own contribution. Every later stage -- memoization,
tabulation, space optimization -- carries over with zero additional
changes, exactly as it did for the 1D Fibonacci-to-Frog-Jump
transformation earlier in this book.
\end{takeawaybox}
